% moi_theorems: the parallel- and perpendicular-axis theorems (Ch 14, fig:theorems)
% Left (parallel axis, any body): a deliberately irregular lamina-like outline drawn in
%   the page, r(phi) = 1.35 (1 + 0.16 cos phi + 0.10 sin 2phi - 0.06 cos 3phi).  For a
%   uniform plate of this shape the centroid (computed numerically) is at
%   (0.214, 0.029); the CM axis passes through it and a parallel axis lies d = 1.74 away.
%   (The theorem holds for any body; the outline only has to have its CM marked right.)
% Right (perpendicular axis, laminae only): a flat plate seen obliquely, with x and y in
%   its plane and z perpendicular to it, all through one point O of the plate.
% Rotation axes use the dash-dot style of the moid_* figures.
\begin{tikzpicture}[x=1cm,y=1cm]
  % ------------------------------------------------------------ parallel axis
  \def\xc{0.214} \def\yc{0.029}   % centre of mass
  \def\xd{1.95}                   % the parallel axis
  \draw[body,fill=cBlock!35] plot[domain=0:360,samples=180,smooth cycle,variable=\p]
      ({1.35*(1+0.16*cos(\p)+0.10*sin(2*\p)-0.06*cos(3*\p))*cos(\p)},
       {1.35*(1+0.16*cos(\p)+0.10*sin(2*\p)-0.06*cos(3*\p))*sin(\p)});
  \draw[rotation axis] (\xc,-1.95) -- (\xc,1.85) node[anchor=south] {$I_{\text{cm}}$};
  \draw[rotation axis] (\xd,-1.95) -- (\xd,1.85) node[anchor=south] {$I$};
  \pic at (\xc,\yc) {cm};
  \node[anchor=south east,inner sep=3pt] at (\xc,\yc) {CM};
  \draw[dim] (\xc,-1.72) -- node[anchor=north,cSteel] {$d$} (\xd,-1.72);
  \node[anchor=north,font=\footnotesize,text=cSteel,align=center] at ({(\xc+\xd)/2-0.35},-2.25)
      {any body, any parallel axis};
  % ------------------------------------------------------- perpendicular axis
  % oblique axes: x towards the reader (lower left), y to the right, z up
  \begin{scope}[shift={(6.9,0)},x={(-0.6cm,-0.42cm)},y={(1cm,0cm)},z={(0cm,1cm)}]
    % the plate: |x| <= 1.9, |y| <= 1.6, z = 0
    \draw[body,fill=cBlock!35] (1.9,-1.6,0) -- (1.9,1.6,0) -- (-1.9,1.6,0) -- (-1.9,-1.6,0) -- cycle;
    % the axes through O (z only above the plate, which hides the rest)
    \draw[rotation axis] (-2.9,0,0) -- (2.9,0,0) node[anchor=north east,inner sep=1pt] {$x$};
    \draw[rotation axis] (0,-2.25,0) -- (0,2.25,0) node[anchor=west,inner sep=2pt] {$y$};
    \draw[rotation axis] (0,0,0) -- (0,0,1.85) node[anchor=south] {$z$};
    \fill[cInk] (0,0,0) circle[radius=0.045cm];
    \node[anchor=north west,inner sep=1pt] at (0,0.06,-0.08) {$O$};
    \node[anchor=north,font=\footnotesize,text=cSteel] at (0,0.3,-2.25)
        {a lamina: all the mass in the plane of $x$ and $y$};
  \end{scope}
\end{tikzpicture}
